% !TeX root = ../report.tex
% Packages, fonts, margins, captions, listings and macros (moved from report.tex)
%%%%%%%%%%%%%%%%%%%%%% Packages and Configs %%%%%%%%%%%%%%%%%%%%%%%%%

\usepackage{color}
\usepackage{pgfgantt}                   % Gantt chart
\usepackage{polyglossia}
\usepackage{caption}                    % Figures and tables captions
\usepackage{pdflscape}                 % Landscape pages for wide diagrams

\usepackage{etoolbox}                   % Patching commands for macros creating i.e. preto
\usepackage{tabularx}
\usepackage{longtable}                  % Long continuous table
\usepackage{ragged2e}                   % Forcing figure
\usepackage{float}                      % ER and checkmark
\usepackage{tikz}                       % Drawing graphs and symbols
\usetikzlibrary{positioning, arrows.meta, fit, backgrounds}
\usepackage{subfig}                     % Subfigures
\usepackage{keyval}                     % Subfigures alignment
\usepackage[export]{adjustbox}          % Subfigures aLignment
\newcommand{\framedgraphic}[2][]{\adjustbox{max width=\linewidth}{\fbox{\includegraphics[#1]{#2}}}}
\usepackage{cleveref}                   % Clever label references
\usepackage{xltabular}   % longtable + tabularx X-columns, breaks across pages

\usepackage{lmodern}
\usepackage{listings}                   % Listings as code blocks
\usepackage{packages/listings-golang}   % Listings support for Golang
\usepackage{packages/listings-yaml}     % Listing support for YAML
\usepackage{packages/report-compat}
\usepackage{hhline}

%: Biblatex for managing bibliography; instead of using old \bibliography{}
%: NOTE: Biblatex define typesetting onto compilation cache.
%: So recompile from scratch (Remove aux and bbl file and recompile)
%: Everytime you stopped using biblatex

\usepackage[
  backend=biber,      % Biber or Bibtex
  style=ieee,         % Bib standard
  sorting=none,       % Sort by appearance in paper
  % sorting=ynt,       % Sort by year, name...
  % sortcites=true,    % Some other  example options ...
  dateabbrev=true,
  urldate=long,
  block=none,
  indexing=false,
  citereset=none,
  isbn=true,
  url=true,
  doi=true,           % Prints DOI
  natbib=true         % if you need Natbib functions
]{biblatex}

%%%%%%%%%%%%%%%%%% Custom Macros & Settings %%%%%%%%%%%%%%%%%%%%%%%
%: From this part to \begin{document}
%: Do not modify unless you know what you are doing...

%%%%%%%%%%%%%%%%%%%%%%%%%%%%%%%%%
%%%%%%% Biblatex macro %%%%%%%%%%

%: Setup URL access date string
\DefineBibliographyStrings{english}{%
  urlseen = {accessed}
}

%: Reference bib files
%: Make sure you reference correct bib file...
\addbibresource{bib/cyc.bib}
\DeclareLanguageMapping{thai}{english}

%%%%%%%%%%%%%%%%%%%%%%%%%%%%%%%%%
%%%%%%%%% Font settings %%%%%%%%%

%: Define 1st and 2nd language
\setdefaultlanguage{thai}
\setotherlanguage{english}

%: Uncomment following to used Laksaman
%: which is equivalent of Sarabun
% \setmainfont[
%     Scale=1.0,
%     LetterSpace=0,
%     WordSpace=1.0,
%     FakeStretch=1.0,
%     Mapping=tex-text
% ]{Laksaman:script=thai}
% \newfontfamily\thaifont[
%     Scale=1.0,
%     LetterSpace=0,
%     WordSpace=1.0,
%     FakeStretch=1.0,
%     Mapping=tex-text
% ]{Laksaman:script=thai}
% \newfontfamily\thaifonttt{Tlwg Typist:script=thai}

%: Uncomment to use actual Sarabun (font file required)
\newfontfamily\thaifont[Script=Thai,Scale=1.23, Mapping=tex-text]{THSarabunNew.ttf}[
  Path=fonts/,
  Extension=.ttf,
  BoldFont=*-Bold,
  ItalicFont=*-Italic,
  BoldItalicFont=*-BoldItalic,
]
\newfontfamily\thaifonttt[Script=Thai,Scale=1.23]{THSarabunNew.ttf}[
  Path=fonts/,
  Extension=.ttf,
  BoldFont=*-Bold,
  ItalicFont=*-Italic,
  BoldItalicFont=*-BoldItalic,
]

\defaultfontfeatures{Mapping=tex-text,Scale=1.23,LetterSpace=0.0}
\setmainfont[Scale=1.23,LetterSpace=0,WordSpace=1.0,FakeStretch=1.0,Mapping=tex-text]{THSarabunNew.ttf}[
  Path=fonts/,
  Extension=.ttf,
  BoldFont=*-Bold,
  ItalicFont=*-Italic,
  BoldItalicFont=*-BoldItalic,
]

%: Set linebreak to Thai
\XeTeXlinebreaklocale "th"
\XeTeXlinebreakskip = 0pt plus 1pt minus 0.1pt
\emergencystretch=3em

%: Define fonts for code blocks
\newfontfamily\codefont[Scale=0.9]{CourierPrime.ttf}[
  Path=fonts/,
  Extension=.ttf,
  BoldFont=*-Bold,
  ItalicFont=*-Italic,
  BoldItalicFont=*-BoldItalic,
]

%%%%%%%%%%%%%%%%%%%%%%%%%%%%%%%%%
%%%%%%%% Margin settings %%%%%%%%
%\linespread{1.6} % Double-spaced
%\oddsidemargin    0.5in
%\evensidemargin   0.5in

%%%%%%%%%%%%%%%%%%%%%%%%%%%%%%%%%
%%%%%%%% Content font size %%%%%%%
%: Body text size in real (rendered) points.
%: TH Sarabun New is loaded with Scale=1.23, so the size LaTeX is asked for
%: is divided by 1.23 to make the rendered size come out at exactly this value.
%: Change \bodyfontsize to resize the whole document (headings follow).
\newdimen\bodyfontsize
\bodyfontsize=14.3pt   % matches the progress report example (was 16pt)

\makeatletter
%: \thsz{\cmd}{size*10}{baseline*10} sets \cmd to a fixed \fontsize. The two numbers
%: are the size10.clo values (in pt) times 10. The numbers are computed once here
%: (not at every use). They are declared robust like the ones in size10.clo, which
%: keeps them safe inside \MakeUppercase (used for chapter headings).
\newcommand\thsz[3]{%
  \begingroup
  \edef\thsz@x{\endgroup\noexpand\DeclareRobustCommand\noexpand#1{\noexpand\@setfontsize\noexpand#1%
  {\strip@pt\dimexpr\bodyfontsize*#2/123\relax}{\strip@pt\dimexpr\bodyfontsize*#3/123\relax}}}%
\thsz@x}
\thsz\normalsize{100}{120}
\thsz\tiny{50}{60}
\thsz\scriptsize{70}{80}
\thsz\footnotesize{80}{95}
\thsz\small{90}{110}
\thsz\large{120}{140}
\thsz\Large{144}{180}
\thsz\LARGE{173}{220}
\thsz\huge{207}{250}
\thsz\Huge{249}{300}
%: Display math / list spacing that size10.clo attaches to these commands
\g@addto@macro\normalsize{%
  \abovedisplayskip 12\p@ \@plus3\p@ \@minus 6\p@
  \abovedisplayshortskip \z@ \@plus 3\p@
  \belowdisplayshortskip 6.5\p@ \@plus3.5\p@ \@minus 3\p@
  \belowdisplayskip \abovedisplayskip
\let\@listi\@listI}
\g@addto@macro\footnotesize{%
  \abovedisplayskip 6\p@ \@plus2\p@ \@minus4\p@ \abovedisplayshortskip \z@ \@plus\p@
\belowdisplayshortskip 3\p@ \@plus\p@ \@minus2\p@ \belowdisplayskip \abovedisplayskip}
\g@addto@macro\small{%
  \abovedisplayskip 8.5\p@ \@plus3\p@ \@minus4\p@ \abovedisplayshortskip \z@ \@plus2\p@
\belowdisplayshortskip 4\p@ \@plus2\p@ \@minus2\p@ \belowdisplayskip \abovedisplayskip}

%: The class hard-codes heading sizes (12pt / 11pt) for a 10pt body; with a 16pt
%: body they would render smaller than the text. Same layout as cpe-thesis-th.cls,
%: but the sizes keep the class' original ratios to the body (1.2x / 1.1x).
\thsz\thsz@section{120}{168}
\thsz\thsz@subsection{110}{156}
\renewcommand
\section{\@startsection{section}{1}{\z@}%
  {-3.5ex \@plus -1ex \@minus -.2ex}{0.1ex\@plus.1ex}%
{\normalfont\thsz@section\bfseries}}
\renewcommand
\subsection{\@startsection{subsection}{2}{\z@}%
  {-3.25ex\@plus -1ex \@minus -.2ex}{0.1ex \@plus .2ex}%
{\normalfont\thsz@subsection\bfseries}}
\makeatother

%%%%%%%%%%%%%%%%%%%%%%%%%%%%%%%%%
%%%%%% Paragraph Settings %%%%%%%
\linespread{1.15}

\newcommand{\thaijustify}[1]{%
  % \ifnum\value{isfigure}=0
  %   \newline % If counter is set, no newline
  % \fi
  \par\hspace{30pt}\justifying
  #1
  % \vspace{1mm}
  % \setcounter{isfigure}{0}%
}

\newcommand{\thaicaption}[1]{%
  % \ifnum\value{isfigure}=0
  %   \newline % If counter is set, no newline
  % \fi
  \par\justifying
  #1
  % \vspace{1mm}
  % \setcounter{isfigure}{0}%
}

%%%%%%%%%%%%%%%%%%%%%%%%%%%%%%%%%
%%%%%% Caption settings %%%%%%%%%
%: Define how the captions of each report
%: element are formatted.

\captionsetup[table]{
  format=hang,
  labelfont=bf,
  position=top,
  labelsep=space,
  justification=justified,
  singlelinecheck=true
}
\captionsetup[figure]{
  format=hang,
  labelfont=bf,
  position=bottom,
  labelsep=space,
  justification=justified,
  singlelinecheck=true
}
\captionsetup[subfloat]{
  format=hang,
  labelfont=bf,
  labelformat=simple,
  labelsep=period,
  justification=justified,
  singlelinecheck=true
  %listofformat=subparens
}
\captionsetup[lstlisting]{
  format=hang,
  labelfont=bf,
  position=top,
  labelsep=space,
  justification=justified,
  singlelinecheck=true,
  skip=10pt
}
%%%%%%%%%%%%%%%%%%%%%%%%%%%%%%%%%
%%%%%%% Numbering styling %%%%%%%
% Enumerate using arabic numbers
\renewcommand{\labelenumii}{\arabic{enumi}.\arabic{enumii}}
\renewcommand{\labelenumiii}{\arabic{enumi}.\arabic{enumii}.\arabic{enumiii}}
\renewcommand{\labelenumiv}{\arabic{enumi}.\arabic{enumii}.\arabic{enumiii}.\arabic{enumiv}}

%%%%%%%%%%%%%%%%%%%%%%%%%%%%%%%%%
%%%%%% Subfigure settings %%%%%%%

%\renewcommand{\subfigtopskip}{10pt}
%\renewcommand{\subfigbottomskip}{-5pt}
%\renewcommand{\subfigcapskip}{-6pt} %vertical space between caption
%                                    %and figure.
%\renewcommand{\subfigcapmargin}{0pt}

% Use thai alphabets for captions
\renewcommand{\thesubfigure}{\thalph{subfigure}}

%%%%%%%%%%%%%%%%%%%%%%%%%%%%%%%%
%%%%%%% Tikz Settings %%%%%%%%%%
%: Tikz is used for drawing simple edge and node graph
%: Import shapes for usage
\usetikzlibrary{er, positioning, shapes.geometric, arrows}

%: Define checkmark symbol
\def\checkmark{\tikz\fill[scale=0.4](0,.35) -- (.25,0) -- (1,.7) -- (.25,.15) -- cycle;}

%%%%%%%%%%%%%%%%%%%%%%%%%%%%%%%%
%%%%%% Listings Settings %%%%%%%%
%: Listings is used for displaying the program or
%: codes in reports as statement

%: Define text color for code blocks
\definecolor{dkgreen}{rgb}{0,0.6,0}
\definecolor{gray}{rgb}{0.5,0.5,0.5}
\definecolor{mauve}{rgb}{0.58,0,0.82}

%: Define code blocks settings
\lstset{
  frame=tb,
  language=Golang,
  aboveskip=3mm,
  belowskip=3mm,
  showstringspaces=false,
  columns=flexible,
  basicstyle={\small\codefont},
  numbers=left,
  numberstyle=\small\color{gray},
  keywordstyle=\color{blue},
  commentstyle=\color{dkgreen},
  stringstyle=\color{mauve},
  breaklines=false,
  breakatwhitespace=false,
  tabsize=3
}

%%%%%%%%%%%%%%%%%%%%%%%%%%%%%%%%%
%%%%%%%%%% Math macro %%%%%%%%%%%
%: Fraction settings
\renewcommand{\topfraction}{0.85}
\renewcommand{\textfraction}{0.1}

%: Define theorem and proof
\newtheorem{theorem}{Theorem}
\newtheorem{lemma}{Lemma}
\newtheorem{corollary}{Corollary}

%: `proof` and `definition` environments already exist in cpe-thesis-th.cls,
%: so they are intentionally NOT redefined here (a local \def\proof would shadow them).

%\newenvironment{proof}{{\sc Proof:}}{~\hfill \blacksquare}
%% The hyperref package redefines the \appendix. This one
%% is from the dissertation.cls
%\def\appendix#1{\iffirstappendix \appendixcover \firstappendixfalse \fi \chapter{#1}}
%\renewcommand{\arraystretch}{0.8}
%%%%%%%%%%%%%%%%%%%%%%%%%%%%%%%%%
